% ═══════════════════════════════════════════════════════════════════════════════
%  CAPÍTULO — Prototype
% ═══════════════════════════════════════════════════════════════════════════════

\chapter{Prototype}

\patroninfo{Creacional}{117}

% ─── Situación Problema ──────────────────────────────────────────────────────
\section{Situación Problema}

Muchos anfitriones de Airbnb administran más de una propiedad y con
frecuencia desean publicar alojamientos muy similares entre sí: dos
apartamentos en el mismo edificio, o varias habitaciones privadas con
las mismas reglas de la casa, los mismos precios base y las mismas
amenidades. Actualmente, para registrar una segunda propiedad similar,
el anfitrión debe rellenar manualmente todos los campos del formulario
desde cero, repitiendo información que ya existe en la primera.

Copiar los datos de una propiedad a mano es tedioso, propenso a
inconsistencias y frustrante para el anfitrión. Una solución trivial
como compartir la misma referencia de objeto entre dos propiedades es
peligrosa: modificar una afectaría a la otra. Se necesita una copia
independiente que pueda personalizarse sin riesgo.

% ─── Solución ────────────────────────────────────────────────────────────────
\section{Solución}

El patrón Prototype dota a la clase \texttt{Propiedad} de un método
\texttt{clonar()} que produce una copia profunda e independiente del
objeto. El anfitrión selecciona una propiedad existente como plantilla
y solicita duplicarla; el sistema crea un nuevo objeto con todos los
atributos copiados —incluyendo las listas de amenidades y las reglas
de la casa— sobre el cual el anfitrión solo necesita ajustar los campos
que difieren, como el número de habitación o el título del anuncio.

Esto agiliza considerablemente la publicación de propiedades similares,
garantiza que la copia sea totalmente independiente del original y evita
duplicar la lógica de inicialización. El anfitrión obtiene una base
sólida sobre la que trabajar en lugar de un formulario vacío.

% ─── Diagrama de clases ─────────────────────────────────────────────────────
\begin{figure}[H]
    \centering
    % \includegraphics[width=\textwidth]{imagenes/abstract_factory_diagrama.png}
    \caption{Diagrama de clases del patrón Abstract Factory}
\end{figure}


% ─── Roles de las Clases ────────────────────────────────────────────────────
\section{Roles de las Clases}

% Explicar la responsabilidad de cada clase/interfaz

\begin{itemize}
    \item \textbf{AbstractFactory:}
    \item \textbf{ConcreteFactory:}
    \item \textbf{AbstractProduct:}
    \item \textbf{ConcreteProduct:}
    \item \textbf{Client:}
\end{itemize}


% ─── Main ───────────────────────────────────────────────────────────────────
\section{Main}

% Explicar qué realiza el programa principal

\begin{lstlisting}[language=Java, caption=NombreClase]
 
\end{lstlisting}


% ─── Salida del Main ────────────────────────────────────────────────────────
\section{Salida del Main}

\begin{lstlisting}[language=Java, caption=NombreClase]
 
\end{lstlisting}


% ─── Clases ─────────────────────────────────────────────────────────────────
\section{Clases}

% ─── Clase 1 ────────────────────────────────────────────────────────────────
\subsection{NombreClase}

\begin{lstlisting}[language=Java, caption=NombreClase]
 
\end{lstlisting}


% ─── Clase 2 ────────────────────────────────────────────────────────────────
\subsection{NombreClase}

\begin{lstlisting}[language=Java, caption=NombreClase]
 
\end{lstlisting}


% ─── Clase 3 ────────────────────────────────────────────────────────────────
\subsection{NombreClase}

\begin{lstlisting}[language=Java, caption=NombreClase]
 
\end{lstlisting}


% ─── Conclusiones ───────────────────────────────────────────────────────────
\section{Conclusiones}

% Explicar:
% - Ventajas del patrón
% - Cuándo usarlo
% - Beneficios obtenidos
% - Posibles desventajas
% - Relación con principios SOLID